\documentclass[12pt, a4paper]{article} % Classe: article, report o book [8]

% --- Lingua e Codifica (Supporto Italiano) ---
\usepackage[utf8]{inputenc}  % Codifica input
\usepackage[T1]{fontenc}      % Codifica font
\usepackage[italian]{babel}   % Lingua italiana [2]

% --- Impaginazione e Margini ---
\usepackage[margin=2.5cm]{geometry} % Imposta margini (2.5cm su tutti i lati)

% --- Pacchetti Scientifici/Matematici ---
\usepackage{amsmath, amssymb, amsfonts, amsthm} % Matematica avanzata [2]
\usepackage{graphicx} % Per inserire immagini
\usepackage{hyperref} % Per collegamenti ipertestuali
\usepackage{booktabs} % Per tabelle professionali

% --- Formattazione Testo e Altro ---
\usepackage[document]{ragged2e}
\usepackage{xcolor}    % Per usare i colori 

\title{Andrea Gorini}
\date{02-02-2026}

\begin{document}

\maketitle

\section{Ruolo}
Andrea Gorini è un commerciale con esperienza di circa 3 anni all'interno
dell'azienda e si occupa di :
\begin{itemize}
    \item Creazione delle offerte all'interno di Access per conto di altri commerciali.
    \item Creazione di preventivi di commesse , ampliamenti o assistenze.
\end{itemize}

\section{Operazioni comuni}
Le operazioni più comuni svolte da Andrea all'interno dell'applicativo
includono:
\begin{itemize}
    \item Creazione delle offerte per nuovi impianti.
    \item Creazione delle offerte per ampliamenti di impianti esistenti.
    \item Creazione degli ordini qualvolta volta un'offerta viene accettata.
    \item Nella creazione delle offerte ricerca e copia di vecchie offerte simili.
\end{itemize}
\section{Situazione Attuale}

\subsection{Flusso di lavoro attuale:}
L'attuale flusso di lavoro per la creazione delle offerte prevede:
\begin{enumerate}
    \item Creazione di uno schema fuzionale in AutoCAD per definire i componenti
          necessari.
    \item Ricerca di un'offerta simile all'interno del database Access.
    \item Copia dell'offerta esistente e modifica dei dettagli specifici per il nuovo
          cliente.
    \item Aggiunta di componenti specifici basati sullo schema AutoCAD.
    \item Revisione e stampa dell'offerta per l'invio al cliente.
    \item Conversione dell'offerta in ordine una volta accettata dal cliente.
    \item Stampa dei documenti per Ufficio Tecnico e Amministrazione.
\end{enumerate}

\subsection{Criticità dell'attuale sistema}
\begin{itemize}
    \item \textbf{Instabilità:} Il software si blocca frequentemente , specialmente durante la stampa o l'aggiunta di righe, costringendo al riavvio.
    \item \textbf{Mancanza di standardizzazione:} Molte descrizioni \textit{(es. condizioni di pagamento, montaggio)} vengono copiate e incollate da file Word esterni,
          aumentando il rischio di errori o disallineamenti.
    \item \textbf{Codifiche obsolete:} I codici attuali "non parlano" con il resto dell'azienda (Fusion/Ad Hoc) e i nomi dei componenti necessitano di una revisione completa.
    \item \textbf{Aggiornamento prezzi:} La funzione "Aggiorna Prezzi" esiste ma è rischiosa per voci a prezzo zero o componenti speciali \textit{(es. cisterne su misura)}
          che richiedono quotazioni manuali esterne.
\end{itemize}

\subsection{Funzionalità da mantenere}
Andrea sottolinea l'importanza di mantenere alcune funzionalità chiave:
\begin{itemize}
    \item \textbf{Divisione in capito:} La suddivisione dell'offerta in capitoli (es. "Dosaggio", "Quadri Elettrici") è fondamentale per l'organizzazione e la chiarezza.
    \item \textbf{Gestione sconti:} L'attuale sistema di gestione degli sconti è efficace e deve essere preservato.
    \item \textbf{Personalizzazione:} La possibilità di modificare manualmente sconti e condizioni è essenziale per adattarsi alle esigenze specifiche di ogni cliente.
    \item \textbf{Ricerca offerte:} La funzione di ricerca e copia di offerte esistenti è molto utile per velocizzare il processo di creazione.
    \item \textbf{Stampe contratti:} L'attaule contratto ha tutte le informazioni necessarie e deve essere mantenuto.
\end{itemize}

\section{Requisiti fuzionali emersi}

\subsection{Funzionalità principali}
Nell'intervista sono emersi i seguenti requisiti funzionali chiave:
\begin{itemize}
    \item \textbf{Modalità offline:} Requisito critico. Il commerciale deve poter lavorare in assenza di connessione \textit{(es. scantinati, aerei)} con sincronizzazione automatica al ripristino della rete.
    \item \textbf{Versioning delle offerte:} Necessità di gestire le revisioni (R1, R2, R3) mantenendo invariato il numero di offerta principale per migliorare la rintracciabilità del progetto. Un cliente potrebbe avere in trattaiva più ampliamenti e
          se per ogni ampliamento si hanno numri di offerta diversi si ha una porpagazione di numeri difficile da gestire.
    \item \textbf{Inserimento a blocchi/capitoli:} Possibilità di inserire nuovi capitoli o blocchi di componenti tra voci già esistenti (es. inserire un dosaggio tra la voce 4 e la 5) senza dover rifare l'offerta.
    \item \textbf{Gestione prezzi flessibile:} Mantenere l'aggiornamento massivo dei prezzi qual'ora si renda necessaria , ma permettere l'editabilità manuale (sconti strategici, validità offerta, accordi specifici).
    \item \textbf{Possibilità di aggiugere capitoli in modo ordinato:} Attualmente tutte le aggiunte sono fatte in coda all'offerta ,
          Andrea suggerisce di poter scegliere in che punto dell'offerta inserire un nuovo capitolo.
\end{itemize} 

\subsection{User experience (UX) e guidata}
Per supportare i profili Junior e ridurre gli errori:
\begin{itemize}
    \item \textbf{Macro-categorie guidate:} Implementazione di "wizard" di selezione. Esempio: Selezionando "Globo", il sistema deve proporre in automatico gli accessori correlati
          \textit{(es. filtri, portelli antiscoppio, ecc.)} per evitare dimenticanze (che causano perdite economiche).
    \item \textbf{Menu a tendina standard:} Sostituzione del copia-incolla con tendine precompilate per condizioni commerciali, montaggi , pagamneti ecc.
          Mantenere la possibilità di editarli manualmente se necessario.
    \item \textbf{Ricerca per tag:} Utile l'inserimeto di tag per le ricerche , Andrea suggerisce però di rimanere "semplici" per evitare di complicare ulteriormente la ricerca.
    \item \textbf{Importazione dati cliente:} Mantenere la possibilità di importare le codizioni di vendita direttamente dall'anagrafica clienti per evitare ridondanze.
          Serve mantenere la possibilità di editarle manualmente se necessario.
    \item \textbf{Template offerte:} Andrea sconsiglia l'addozione di template rigidi per le offerte in quanto ogni cliente ha esigenze diverse.
          Meglio mantenere un'offerta flessibile e personalizzabile, ma guidando l'utente con le funzionalità sopra descritte.
\end{itemize}

\subsection{Integrazione e flussi aziendali}
La nuova soluzione deve integrarsi con i reparti interni:
\begin{itemize}
    \item \textbf{Codifica univoca:} Il nuovo software deve utilizzare codici allineati con l'ERP (Ad Hoc) e l'Ufficio Tecnico (Fusion/Solidworks).
    \item \textbf{Evoluzione tecnica post-vendita:} L'Ufficio Tecnico (PM/PE) dovrebbe poter accedere all'offerta trasformata in ordine per apportare modifiche tecniche (distinta tecnica vs distinta commerciale) e tracciare le variazioni di costo/componenti.
    \item \textbf{Stampe differenziate:} Generazione automatica di documenti distinti per Amministrazione e Ufficio Tecnico senza dover ripetere la procedura di stampa.
          Attuamente il commerciale deve stampare il contratto in due modalità differenti a seconda del reparto di destinazione.
\end{itemize}

\section{Sviluppi futuri}
Andrea identifica due possibili sviluppi futuri che al momento non sono
prioritari:
\begin{itemize}
    \item \textbf{Integrazione calcoli Excel:} Possibilità futura di lanciare fogli di calcolo (es. dimensionamento quadri elettrici) direttamente dal software e importarne il risultato.
    \item \textbf{Integrazione disegni (CAD):} Possibilità di legare l'offerta a schemi o layout semplificati, ma considerata una funzione avanzata per fasi successive.
\end{itemize}

\section{Strategia di rilascio}
L'intervistato consiglia vivamente di:
\begin{enumerate}
    \item Rilasciare una prima versione che replichi le funzioni attuali ma con maggiore
          stabilità e UX migliorata ("Partire Basic").
    \item Far testare il software a utenti "propositivi" (es. Thomas e Andrea stesso)
          prima del rilascio globale.
    \item Implementare le funzioni avanzate (integrazione Excel/CAD) solo dopo che la
          base è consolidata.
\end{enumerate}

\end{document}